%%%%%%%%%%%%%%%%%%%%%%%%%%%%%%%%%%%%%%%%%%%%%%%%%%%%%%%%%%%%%%%%%%%%%%%%%%
%% The Global South AI Safety Hackathon (June 2026) -- submission template
%% Single-column arXiv-style preprint. Build with: bash build.sh
%%
%% HOW TO WRITE IN THIS TEMPLATE
%%   - Guidance text is wrapped in \guide{...} and renders gray+italic.
%%     Replace each \guide{...} block with your own content, then remove it.
%%   - The yellow info box and all \guide{...} text MUST be deleted before
%%     submitting (see the box on page 1).
%%   - Recommended length: 4 pages excluding references and appendix.
%%%%%%%%%%%%%%%%%%%%%%%%%%%%%%%%%%%%%%%%%%%%%%%%%%%%%%%%%%%%%%%%%%%%%%%%%%

\documentclass[11pt]{article}

% --- Page geometry ---
\usepackage[letterpaper,margin=1in]{geometry}

% --- Encoding ---
\usepackage[T1]{fontenc}
\usepackage[utf8]{inputenc}

% --- Math (amsmath must load before the newtx fonts) ---
\usepackage{amsmath,mathtools}

% --- Fonts: Times-like serif to match the official template.
%     newtxmath supplies the AMS/blackboard symbols, so amssymb is not loaded
%     separately (doing so clashes on \Bbbk). ---
\usepackage{newtxtext}
\usepackage{newtxmath}

% --- Graphics, color, tables ---
\usepackage{graphicx}
\graphicspath{{images/}}
\usepackage{xcolor}
\usepackage{booktabs}
\usepackage{array}
\usepackage{multirow}

% --- Captions / floats ---
\usepackage{caption}
\usepackage{subcaption}
\usepackage{float}

% --- Lists ---
\usepackage{enumitem}

% --- Info box (deletable guidance box) ---
\usepackage[most]{tcolorbox}

% --- Typography ---
\usepackage{microtype}

% --- Section headings: "1. Introduction" (period after the number) ---
\usepackage{titlesec}
\titlelabel{\thetitle.\quad}

% --- Bibliography ---
\usepackage[numbers,sort&compress]{natbib}

% --- Hyperlinks (load late) ---
\usepackage{url}
\usepackage[colorlinks=true,allcolors=blue!55!black]{hyperref}
\usepackage[capitalize,noabbrev]{cleveref}

%%%%%%%%%%%%%%%%%%%%%%%%%%%%%%%%%%%%%%%%%%%%%%%%%%%%%%%%%%%%%%%%%%%%%%%%%%
%% Helper: \guide{...} marks replaceable guidance text (gray italic).
%% Delete every \guide{...} block before submitting.
%%%%%%%%%%%%%%%%%%%%%%%%%%%%%%%%%%%%%%%%%%%%%%%%%%%%%%%%%%%%%%%%%%%%%%%%%%
\definecolor{guidecolor}{RGB}{95,95,95}
\newcommand{\guide}[1]{\par{\color{guidecolor}\itshape #1}\par\medskip}

\begin{document}
\thispagestyle{plain}

%%%%%%%%%%%%%%%%%%%%%%%%%%%%%%%%%%%%%%%%%%%%%%%%%%%%%%%%%%%%%%%%%%%%%%%%%%
%% TITLE BLOCK
%%%%%%%%%%%%%%%%%%%%%%%%%%%%%%%%%%%%%%%%%%%%%%%%%%%%%%%%%%%%%%%%%%%%%%%%%%
\noindent\rule{\textwidth}{1.2pt}\par
\vspace{0.8em}
\begin{center}
  {\LARGE\bfseries PROJECT TITLE%
    \footnote{Research conducted at the \href{https://www.apartresearch.com/}{The Global South AI Safety Hackathon}, June 2026.}\par}
\end{center}
\vspace{0.6em}
\noindent\rule{\textwidth}{0.4pt}\par
\vspace{1.4em}

%% Authors -- 3 across, 2 rows (delete unused cells). Keep affiliations short.
\begin{center}
  \renewcommand{\arraystretch}{1.15}
  \begin{tabular}{ccc}
    Author name 1 & Author name 2 & Author name 3 \\
    \textit{Affiliation} & \textit{Affiliation} & \textit{Affiliation} \\[1.0em]
    Author name 4 & Author name 5 & Author name 6 \\
    \textit{Affiliation} & \textit{Affiliation} & \textit{Affiliation} \\
  \end{tabular}

  \vspace{1.1em}
  {\bfseries With}\\[0.2em]
  Apart Research
\end{center}
\vspace{1.2em}

%%%%%%%%%%%%%%%%%%%%%%%%%%%%%%%%%%%%%%%%%%%%%%%%%%%%%%%%%%%%%%%%%%%%%%%%%%
%% ABSTRACT
%%%%%%%%%%%%%%%%%%%%%%%%%%%%%%%%%%%%%%%%%%%%%%%%%%%%%%%%%%%%%%%%%%%%%%%%%%
\begin{abstract}
\guide{Summarize your project in 150--250 words. A strong abstract lets a
reviewer understand what you did and why it matters without reading anything
else. Make sure to cover: the problem, your approach, key results, and the main
takeaway. Polish it last: the abstract should reflect your final results, not
your initial plan.}
\end{abstract}

\bigskip

%%%%%%%%%%%%%%%%%%%%%%%%%%%%%%%%%%%%%%%%%%%%%%%%%%%%%%%%%%%%%%%%%%%%%%%%%%
%% INFO BOX -- DELETE THIS ENTIRE BOX BEFORE SUBMITTING
%%%%%%%%%%%%%%%%%%%%%%%%%%%%%%%%%%%%%%%%%%%%%%%%%%%%%%%%%%%%%%%%%%%%%%%%%%
\begin{tcolorbox}[colback=yellow!12,colframe=yellow!55!black,boxrule=0.4pt,
  arc=1pt,left=8pt,right=8pt,top=8pt,bottom=8pt]
\textbf{How to use this template}: Replace the gray italic guidance text
(\texttt{\textbackslash guide\{\ldots\}}) under each section with your content.
The section structure is strong guidance but not rigid. If your project
requires a different organization, feel free to adapt. \underline{Delete all
guidance text including this info box before submitting.} Make sure the project
title and author information above are up to date.

\medskip
\textbf{How your submission will be evaluated}: Your project will be judged on
the quality of this written report. To understand what we look for, see the
\href{https://www.apartresearch.com/}{Evaluation Rubric}.

\medskip
\textbf{Recommended length}: 4 pages excluding references and appendix.\\
Rough guide: Intro \& Related Work 1p, Methods and Results 2.5p, Discussion 0.5p.
\end{tcolorbox}

%%%%%%%%%%%%%%%%%%%%%%%%%%%%%%%%%%%%%%%%%%%%%%%%%%%%%%%%%%%%%%%%%%%%%%%%%%
\section{Introduction}
\guide{What problem are you addressing and why does it matter? Connect it to
your work: we want to know why your work is practically valuable.}

\guide{Provide enough background for readers to understand your work. If
relevant, briefly describe the threat model or failure mode you're addressing;
reference prior work that motivates why it is worth addressing, or explain it
yourself.}

\guide{Aspire to clearly list your most important contributions that go beyond
what exists today.}

\noindent\textbf{Our main contributions are:}
\begin{enumerate}[leftmargin=*,topsep=2pt]
  \item \guide{First contribution --- what new thing did you create, discover, or demonstrate?}
  \item \guide{Second contribution.}
  \item \guide{Third contribution, if applicable.}
\end{enumerate}

%%%%%%%%%%%%%%%%%%%%%%%%%%%%%%%%%%%%%%%%%%%%%%%%%%%%%%%%%%%%%%%%%%%%%%%%%%
\section{Related Work}
\guide{What prior work is most similar, and how does your work differ? Cite the
most relevant papers, tools, or projects~\citep{example}. Explain what gap your
work addresses.}

\guide{Some questions which may help:}
\begin{enumerate}[label=(\alph*),leftmargin=*,topsep=2pt]
  \item \guide{When and why would someone use your method over the existing state-of-the-art?}
  \item \guide{What information/insight does your method provide which we did not have before?}
\end{enumerate}

%%%%%%%%%%%%%%%%%%%%%%%%%%%%%%%%%%%%%%%%%%%%%%%%%%%%%%%%%%%%%%%%%%%%%%%%%%
\section{Methods}
\guide{Describe your approach clearly enough that someone could replicate it.
Include key design choices and justify them where relevant (hint: the more you
can back up your design choices by referencing prior work, the better).}

\guide{What models/datasets/tools did you use and why? What were key parameters
or design decisions? What did you try that didn't work? Could someone reproduce
your work from this description?}

%%%%%%%%%%%%%%%%%%%%%%%%%%%%%%%%%%%%%%%%%%%%%%%%%%%%%%%%%%%%%%%%%%%%%%%%%%
\section{Results}
\guide{Present your main findings with appropriate evidence. Use figures and
tables where appropriate (we strongly encourage at least one figure, see tips
below). Distinguish between observations and interpretations.}

\guide{Argue why your claims are robust. E.g., if your approach ``performs
better'' than alternatives:}
\begin{enumerate}[label=(\alph*),leftmargin=*,topsep=2pt]
  \item \guide{Do you have enough data? Is the difference statistically significant?}
  \item \guide{Is it robust? Or do small changes to your setup cause substantial changes to your results?}
\end{enumerate}

\noindent\textbf{Tips for figures and tables:}
\begin{itemize}[label=\textendash,leftmargin=*,topsep=2pt]
  \item \guide{Number all figures (Figure 1, Figure 2\ldots) and tables (Table 1, Table 2\ldots).}
  \item \guide{Include descriptive captions that can be understood without the main text.}
  \item \guide{Place figures/tables near where they're first referenced.}
  \item \guide{Ensure text in figures is legible!}
\end{itemize}

%% Example figure scaffold (uncomment and drop a file in images/):
% \begin{figure}[t]
%   \centering
%   \includegraphics[width=0.8\linewidth]{example.png}
%   \caption{A self-contained caption.}
%   \label{fig:example}
% \end{figure}

%%%%%%%%%%%%%%%%%%%%%%%%%%%%%%%%%%%%%%%%%%%%%%%%%%%%%%%%%%%%%%%%%%%%%%%%%%
\section{Discussion and Limitations}
\guide{Discuss the broader implications for AI safety. What do your results
mean? What trends do you notice and what might they indicate?}

\subsection*{Limitations}
\guide{What are the limitations of your work? What threat models or failure
modes did you not address? Be honest about constraints --- methodological
limitations, scope limitations, or aspects you couldn't fully address in the
hackathon timeframe. Explicitly note the assumptions you made, whether
implicitly or explicitly, and how the interpretation of your results would
change if a given assumption did not hold.}

\subsection*{Future Work}
\guide{What are the natural next steps? How could this work be extended?}

%%%%%%%%%%%%%%%%%%%%%%%%%%%%%%%%%%%%%%%%%%%%%%%%%%%%%%%%%%%%%%%%%%%%%%%%%%
\section{Conclusion}
\guide{Briefly summarize your main findings and their implications (1--2 paragraphs).}

%%%%%%%%%%%%%%%%%%%%%%%%%%%%%%%%%%%%%%%%%%%%%%%%%%%%%%%%%%%%%%%%%%%%%%%%%%
\section*{Code and Data}
\guide{Include links if applicable. If your project doesn't involve code (e.g.,
policy analysis) or if there are info-hazard considerations, note that here.}
\begin{itemize}[label=\textendash,leftmargin=*,topsep=2pt]
  \item \textbf{Code repository:} \guide{Link to GitHub/GitLab if applicable.}
  \item \textbf{Data/Datasets:} \guide{Link if applicable.}
  \item \textbf{Other artifacts (optional):} \guide{Demo link, video walkthrough, Hugging Face Space, etc.}
\end{itemize}

%%%%%%%%%%%%%%%%%%%%%%%%%%%%%%%%%%%%%%%%%%%%%%%%%%%%%%%%%%%%%%%%%%%%%%%%%%
\section*{Author Contributions (optional)}
\guide{e.g., ``A.B. led the project and designed experiments. C.D. implemented
the code. All authors contributed to writing and reviewed the final manuscript.''}

%%%%%%%%%%%%%%%%%%%%%%%%%%%%%%%%%%%%%%%%%%%%%%%%%%%%%%%%%%%%%%%%%%%%%%%%%%
%% REFERENCES -- add entries to references.bib and cite with \citep{key}.
%%%%%%%%%%%%%%%%%%%%%%%%%%%%%%%%%%%%%%%%%%%%%%%%%%%%%%%%%%%%%%%%%%%%%%%%%%
\bibliographystyle{plainnat}
\bibliography{references}

%%%%%%%%%%%%%%%%%%%%%%%%%%%%%%%%%%%%%%%%%%%%%%%%%%%%%%%%%%%%%%%%%%%%%%%%%%
\section*{Appendix (optional)}
\guide{Supplementary material such as additional figures, detailed methodology,
prompts used, extended results, etc.}

%%%%%%%%%%%%%%%%%%%%%%%%%%%%%%%%%%%%%%%%%%%%%%%%%%%%%%%%%%%%%%%%%%%%%%%%%%
\section*{LLM Usage Statement}
\guide{If you used LLM assistance in developing your project or writing this
report, briefly note how. Ensure all claims and results have been verified.
NOTE: We strongly encourage that the final version of the submission is
primarily written by your team. e.g., ``We used Claude to brainstorm approaches
and help draft sections. All results and claims were independently verified.''}

\end{document}
